\documentclass[aspectratio=169]{beamer}

% ===== 基本パッケージ =====
\usepackage{adjustbox}
\usepackage{appendixnumberbeamer}
\usepackage[english]{babel}
\usepackage[utf8x]{inputenc}
\usepackage{amssymb,booktabs,amsfonts,mathtools}
\usepackage[babel]{csquotes}
\usepackage{subfiles}
\usepackage{threeparttable}
\usepackage{graphicx}
\usepackage[round]{natbib}
\usepackage{tabularray}
\usepackage{float}
\usepackage[normalem]{ulem}
\usepackage{multirow}
\usepackage{multicol}
\usepackage{comment}
\usepackage{tikz}
\usetikzlibrary{arrows.meta}

% ===== フォント =====
\usefonttheme{professionalfonts}
\renewcommand{\familydefault}{\sfdefault}

% ===== カラー定義 =====
\definecolor{myOrange}{RGB}{230,120,0}
\definecolor{myGreen}{RGB}{0,130,60}   % new (endline) slides
\definecolor{myDarkGray}{RGB}{50,50,50}

% ===== カラーテーマ設定 =====
\setbeamercolor{structure}{fg=myOrange}
\setbeamercolor{title}{fg=myOrange}
\setbeamercolor{frametitle}{fg=white, bg=myOrange}
\setbeamercolor{normal text}{fg=myDarkGray, bg=white}

% ===== フレームタイトル =====
\setbeamertemplate{frametitle}{
	\vspace{0.2em}
	\begin{beamercolorbox}[wd=\paperwidth,ht=2.8ex,dp=1.2ex,leftskip=1em]{frametitle}
		\insertframetitle
	\end{beamercolorbox}
}

% ===== フッター（左：著者 / 中：スライドタイトル / 右：ページ）=====
\setbeamertemplate{footline}{
	\leavevmode
	\hbox{
		% 左
		\begin{beamercolorbox}[wd=.33\paperwidth,ht=2.5ex,dp=1.2ex,left]{author in head/foot}
			\hspace{1em}\insertshortauthor
		\end{beamercolorbox}
		% 中
 		 \begin{beamercolorbox}[wd=.33\paperwidth,ht=2.5ex,dp=1.2ex,center]{title in head/foot}% \hspace{1em}
			\usebeamerfont{title in head/foot}
		\end{beamercolorbox}%
		% 右
		\begin{beamercolorbox}[wd=.33\paperwidth,ht=2.5ex,dp=1.2ex,right]{date in head/foot}
			\insertframenumber{} / \inserttotalframenumber\hspace{1em}
		\end{beamercolorbox}
	}
}

\setbeamercolor{author in head/foot}{fg=white,bg=myOrange}
\setbeamercolor{title in head/foot}{fg=white,bg=myOrange}
\setbeamercolor{date in head/foot}{fg=white,bg=myOrange}

% ===== ナビゲーション消去 =====
\setbeamertemplate{navigation symbols}{}

% ===== 情報 =====
\makeatletter
\newcommand*\exInput[1]{\@@input#1 }
\makeatother

\title[Short Title]{Results of Baseline, Bi-weekly, and Endline Surveys}
\author[Takahashi et al.]{Kazushi Takahashi, Wataru Kodama, Yukichi Mano}
\date{\today}
\begin{document}
	
	% --- 1枚目のスライド（タイトルページ） ---
	\begin{frame}{Title}
		\titlepage
	\end{frame}
	
	% --- 2枚目のスライド ---
	\begin{frame}{Contents} % 
		\begin{itemize}
			\item Motivation % 
			\item Research Questions and What We Do
			\item Study Areas and Research Design
			\item Descriptive statistics
			\item Preliminary results
			\item Endline impact results
		\end{itemize}
	\end{frame}
% --- Motivation ---
\begin{frame}{Motivation}
	
	\begin{itemize}
		\item Improving agricultural productivity requires the adoption of new technologies,
		such as improved seeds, chemical fertilizers, and appropriate crop management practices.
		
		\vspace{0.3cm}
		
		\item Since technology development, adoption, and diffusion involve externality,
		public institutions play a central role.
		Public agricultural extension workers are key actors in disseminating technologies
		developed at research stations.
		
		\vspace{0.3cm}
		
		\item However, in many developing countries, particularly in Sub-Saharan Africa,
		public extension systems often function poorly.
		
		\vspace{0.3cm}
		
		\item One explanation is weak incentives:
		\begin{itemize}
			\item Limited monitoring and weak linkage between performance and rewards
			(lack of extrinsic motivation)
			\item Extension workers may not perceive extension activities as meaningful
			(lack of intrinsic motivation)
		\end{itemize}
		
		\vspace{0.3cm}
		
		
	\end{itemize}
	
\end{frame}

% --- Research gap \& Aim of the paper ---
\begin{frame}{Research Questions} % 
	\begin{itemize}
		\item This study focuses on intrinsic motivation and evaluates interventions
		that can improve the effectiveness of extension activities.
		
		\vspace{0.3cm}
		
		\begin{itemize}
			
			\item \textbf{Is low extension effort driven by a lack of locus of control?}
			\begin{itemize}
				\item Extension workers do not believe that their efforts can improve farmers' lives.
			\end{itemize}
			
			\vspace{0.3cm}
			
			\item \textbf{Is low effort driven by weak social recognition?}
			\begin{itemize}
				\item No matter how hard AEAs try, their efforts are not particularly recognized or appreciated by their superiors or the farmers. 
			\end{itemize}
			
			\vspace{0.3cm}
			
			
		\end{itemize}
		
	\end{itemize}
\end{frame}

\begin{frame}{What We Do}
	
	\begin{itemize}
		\item This study focuses on intrinsic motivation and evaluates interventions
		that can improve the effectiveness of extension activities.
		
		\vspace{0.3cm}
		\begin{itemize}
			\item \textbf{Is low extension effort driven by a lack of locus of control?}
			\begin{itemize}
				\item Extension workers do not believe that their efforts can improve farmers' lives.
				\item $\Rightarrow$ Provide training emphasizing the social significance of extension activities
			\end{itemize}
			
			\vspace{0.3cm}
			
			\item \textbf{Is low effort driven by weak social recognition?}
			\begin{itemize}
				\item  No matter how hard AEAs try, their efforts are not particularly recognized or appreciated by their superiors or the farmers. 
				\item $\Rightarrow$ Introduce a feedback system communicating farmers' satisfaction
				\item $\Rightarrow$ Ask AEAs to stamp their location information when visit farmer fields (somewhat related to extrinsic motivation?)
			\end{itemize}
			
			\vspace{0.3cm}
			
			
		\end{itemize}
	\end{itemize}
	
\end{frame}

\begin{frame}{Study Areas}
	\begin{itemize}
		\item Rainfed areas covered by GRIP:
		\begin{itemize}
			\item Oti (Biakoye and Krachi West)
			\item Bono East (Pru West and Sene West)
			\item Western North (Bibiani and Sefwi Wiaoso)
			\item Eastern Regions (Achiase and Lower Manya) 
			\vspace{0.3cm}
		\end{itemize}
		\item Irrigation schemes:
		\begin{itemize}
			\item Weta (Volta)
			\item Kpong (Eastern)
		\end{itemize}
		\vspace{0.3cm}
		\item Total: 
		\begin{itemize}
			\item 10 study sites across major rice-producing environments in Ghana
			\item 44 agricultural extension agents (AEAs)
			\item 390 rice-growing farmers
			
		\end{itemize}
		
	\end{itemize}
\end{frame}

\begin{frame}{Random Assignment}
	\textbf{Sample}
	\begin{itemize}
		\item 44 AEAs who participated in the baseline survey
		\item Random assignment conducted at the district level
	\end{itemize}
	
	
	\textbf{Treatment Arms}
	\begin{itemize}
		\item \textbf{Group 1: Feedback Treatment}
		\begin{itemize}
			\item Extension agents and their supervisors receive feedback from farmers regarding their extension activities: biweekly
			\item \textcolor{blue}{Sene West, Achiase, Bibiani}
		\end{itemize}
		
		\item \textbf{Group 2: Feedback + Training}
		\begin{itemize}
			\item Farmer feedback provided
			\item Additional training emphasizing the social importance and impact of extension services
			\item \textcolor{blue}{Lower Manya, Krachi West, Sefwi Wiaoso}
		\end{itemize}
		
		\item \textbf{Group 3: Control Group}
		\begin{itemize}
			\item No intervention (Collect the satisfaction level, but do not inform the result to AEAs) 
			\item \textcolor{blue}{Pru West, Biakoye, and 2 irrigation schemes}
		\end{itemize}
		
		\item \textbf{Half of AEAs are subject to stamp treatment}
		
	\end{itemize}
	
\end{frame}

\begin{frame}{Main outcomes of interest}
	\begin{itemize}
		\item Rice Production (Farmer data at baseline and endline)
		\begin{itemize}
			\item Yield
			\item Income 
			\item Profits
			\item Rice management practices
			\vspace{0.3cm}
		\end{itemize}
		\item Extension Activities (Biweekly data) 
		\begin{itemize}
			\item Number of visits
			\item Number of successful communication
			\item Satisfaction level 
		\end{itemize}
		
		
		
	\end{itemize}
\end{frame}

\begin{frame}{Project Timeline}
	
	\begin{itemize}
		\item \textbf{January 2025:} Baseline Survey
		\vspace{0.3cm}
		
		\item \textbf{April -- September 2025:} Intervention
		\begin{itemize}
			\item Satisfaction Survey
			\item Training
		\end{itemize}
		\vspace{0.3cm}
		
		\item \textbf{January 2026:} Endline Survey
	\end{itemize}
	
\end{frame}

\begin{frame}{Number of Obs by District}
\begin{table}[hbtp]
	\centering
	\begin{threeparttable}
		
		% ここで列の定義を自分で行う
		\begin{tabular}{llccc}
			\hline
			\input{tmp/obs_data.tex} % Stataが出力した中身だけを読み込む
			
		\end{tabular}
		
	\end{threeparttable}
\end{table}
\end{frame}	




\begin{frame}{Demographic Characteristics (Farmer)}
	\centering
	\scalebox{0.55}{\input{tmp/f_balance1.tex}}
	
\end{frame}

\begin{frame}{Basic Farming Characteristics (Farmer)}
	\centering
	\scalebox{0.5}{\input{tmp/f_balance2.tex}}
	
\end{frame}


\begin{frame}{Basic Farming Characteristics (Farmer) w/o Irrg}
\centering
\scalebox{0.5}{\input{tmp/f_balance2b.tex}}

\end{frame}

\begin{frame}{Technology Adoption (Farmer)}
	\centering
	\scalebox{0.6}{\input{tmp/f_balance3.tex}}
	
\end{frame}

\begin{frame}{Basic Characteristics (AEA)}
	\centering
	\scalebox{0.55}{\input{tmp/a_balance1.tex}}
	
\end{frame}

\begin{frame}{Prosocialness (AEA)}
	\centering
	\scalebox{0.6}{\input{tmp/a_balance2.tex}}
	
\end{frame}

\begin{frame}{Intrinsic and Extrinsic Motivation}
	\begin{itemize}
		\item Intrinsic motivation
		\begin{itemize}
			\item Because this activity is fun 
			\item Because I feel good when doing this activity
			\vspace{0.3cm}
		\end{itemize}
		\item Extrinsic Motivation (External regulation and identified regulation)
		\begin{itemize}
			\item Because I am supposed to do it
			\item Because I don't have any choice
			\item Because I am doing it for my own good 
			\item Because I think that this activity is good for me. 
		\end{itemize}
		
		
		
	\end{itemize}
\end{frame}


\begin{frame}{Distribution of Prosocialness Answer}
	\begin{figure}
		\includegraphics[width=12cm]{tmp/prosocial.eps}
	\end{figure}
\end{frame}

\begin{frame}{Locus of Control (AEA)}
	\centering
	\scalebox{0.5}{\input{tmp/a_balance3.tex}}
	
\end{frame}

\begin{frame}{Distribution of Locus of Control Answer (A-E)}
	\begin{figure}
		\includegraphics[width=12cm]{tmp/locus1.eps}
	\end{figure}
\end{frame}

\begin{frame}{Distribution of Locus of Control Answer (F-J)}
	\begin{figure}
		\includegraphics[width=12cm]{tmp/locus2.eps}
	\end{figure}
\end{frame}

\begin{frame}{Descriptive Results of Biweekly Survey}
	\centering
\scalebox{0.6}{\input{tmp/b_balance1.tex}}
\end{frame}




\begin{frame}{Descriptive Results of Biweekly Survey}
	\begin{figure}
		\includegraphics[width=12cm]{tmp/b_basic.eps}
	\end{figure}
\end{frame}

\begin{frame}{Satisfaction Level by District and Time}
	\begin{figure}
		\includegraphics[width=10cm]{tmp/b_satisfy.eps}
	\end{figure}
\end{frame}


\begin{frame}{Preliminary results (AEA)}
	\centering
	\scalebox{0.55}{\input{tmp/result_table1.tex}}
		\begin{tablenotes}
			\scriptsize
		\item Note: Provincial and Round fixed effects are controlled. ***, **, and * indicate significance at the 1, 5, and 10 percent critical level.  \\.  
	\end{tablenotes}
	
\end{frame}

\begin{frame}{Preliminary results (AEA)}
	\centering
	\scalebox{0.55}{\input{tmp/result_table2.tex}}
	\begin{tablenotes}
		\scriptsize
		\item Note: Provincial and Round fixed effects are controlled. ***, **, and * indicate significance at the 1, 5, and 10 percent critical level.  \\.  
	\end{tablenotes}
	
\end{frame}

\begin{frame}{Preliminary results of Correlation (Farming)}
	\centering
	\scalebox{0.34}{\input{tmp/f_result_table2.tex}}
	
\end{frame}

\begin{frame}{Preliminary results of Correlation (AEA)}
	\centering
	\scalebox{0.5}{\input{tmp/a_result_table3.tex}}
	
\end{frame}

% ======================= ENDLINE IMPACT RESULTS =======================
% New slides (Feb 2026 endline): theme switches to GREEN from here so
% co-authors can see at a glance which slides were added.
\setbeamercolor{structure}{fg=myGreen}
\setbeamercolor{frametitle}{fg=white, bg=myGreen}
\setbeamercolor{author in head/foot}{fg=white,bg=myGreen}
\setbeamercolor{title in head/foot}{fg=white,bg=myGreen}
\setbeamercolor{date in head/foot}{fg=white,bg=myGreen}

\begin{frame}{Endline: Sample and Attrition}
\begin{itemize}
	\item Endline survey conducted January--February 2026 (one wet season after the intervention)
	\item Baseline: 390 farmers, 44 AEAs \; $\rightarrow$ \; Endline: 394 farmers, 43 AEAs
	\begin{itemize}
		\item \textbf{385 matched panel households} ($\approx$99\% of baseline)
		\item 1 AEA retired (Worawora); response rates $\approx$98\%
	\end{itemize}
	\item Attrition concentrated in Kpong irrigation and Achiase, largely non-random
	      (plots leased out; farmers exited rice)
	\item Survey-firm data revisions (Jul 2026) incorporated: labor in-kind values,
	      variety/seed-type coding, AEA visit frequencies
\end{itemize}
\end{frame}

\begin{frame}{Descriptives by Treatment Arm (AEA)}
	\centering
	\scalebox{0.6}{\input{tmp/desc_arm_aea.tex}}
	\begin{itemize}\footnotesize
		\item Note: baseline and endline means by assigned arm; SD in parentheses. Stars on T1/T2:
		difference vs Control \emph{within the wave} (district-clustered); ***, **, * : 1, 5, 10\%.
		Self-reported visits winsorized at 1/99 in the analysis (raw here).
	\end{itemize}
\end{frame}

\begin{frame}{Descriptives by Treatment Arm (Farmer)}
	\centering
	\scalebox{0.42}{\input{tmp/desc_arm_farmer.tex}}
	\begin{itemize}\scriptsize
		\item Note: baseline and endline means by assigned arm (baseline cells: matched panel;
		endline cells: endline sample, N=394); SD in parentheses. Stars on T1/T2: difference vs
		Control \emph{within the wave} (district-clustered); ***, **, * : 1, 5, 10\%.
		Income can be negative (crop failure $=0$ output $-$ costs).
	\end{itemize}
\end{frame}

\begin{frame}{Descriptives by Treatment Arm (Farmer, excl.\ Irrigation Schemes in C)}
	\centering
	\scalebox{0.42}{\input{tmp/desc_arm_farmer_noirr.tex}}
	\begin{itemize}\scriptsize
		\item Note: as the previous slide, but the Control group \emph{excludes the two
		irrigation-scheme sites} (Kpong and Weta; 12 endline households). T1/T2 unchanged.
		***, **, * : 1, 5, 10\%.
	\end{itemize}
\end{frame}

\begin{frame}{Prosocialness Answers: Baseline vs Endline (AEA)}
\begin{figure}
	\centering
	\includegraphics[height=0.78\textheight,keepaspectratio]{tmp/prosocial_be.pdf}
\end{figure}
\end{frame}

\begin{frame}{Locus of Control Answers (A--E): Baseline vs Endline (AEA)}
\begin{figure}
	\centering
	\includegraphics[height=0.78\textheight,keepaspectratio]{tmp/locus1_be.pdf}
\end{figure}
\end{frame}

\begin{frame}{Locus of Control Answers (F--J): Baseline vs Endline (AEA)}
\begin{figure}
	\centering
	\includegraphics[height=0.78\textheight,keepaspectratio]{tmp/locus2_be.pdf}
\end{figure}
\end{frame}

\begin{frame}{Descriptives: AEA Outcomes, Baseline vs Endline}
	\centering
	\scalebox{0.75}{\input{tmp/desc_be_aea.tex}}
	\begin{itemize}\footnotesize
		\item Note: matched panel AEAs (N=43). SD in parentheses. Diff = endline $-$ baseline;
		stars from a paired test with district-clustered SEs. ***, **, * : 1, 5, 10\%.
	\end{itemize}
\end{frame}

\begin{frame}{Descriptives: Farmer Outcomes, Baseline vs Endline}
	\centering
	\scalebox{0.52}{\input{tmp/desc_be_farmer.tex}}
	\begin{itemize}\scriptsize
		\item Note: matched panel households (N=385). SD in parentheses. Diff = endline $-$ baseline;
		stars from a paired test with district-clustered SEs.
		``Any improved variety'' pools all improved codes (CSIR-AGRA, Exbaika, Jasmin85,
		Aromatic, Thai Jasmine, Amankwatia). Endline variety/seed-type coding errors were
		corrected by the survey firm (revised data, Jul 2026). ***, **, * : 1, 5, 10\%.
	\end{itemize}
\end{frame}

\begin{frame}{Empirical Strategy}
\begin{itemize}
	\item \textbf{ANCOVA}: regress endline outcome on treatment status, controlling for
	      the baseline value of the outcome
	\item Arms: T1 (Feedback), T2 (Feedback + Training); plus the cross-cutting Stamp (GPS)
	\item Farmer models add region (province) fixed effects; SEs clustered at district ($\approx$10 clusters)
	\item Few clusters $\Rightarrow$ we report \textbf{wild cluster bootstrap} (Webb, 9{,}999) and
	      \textbf{randomization-inference} $p$-values alongside clustered SEs
	\item Crop failure = realized \textbf{zero} output (kept in sample); logs use $\ln(1+x)$;
	      continuous outcomes winsorized at 1/99
\end{itemize}
\end{frame}

\begin{frame}{Bi-weekly (AEA): Effort Level Reported by Farmers}
	\centering
	\scalebox{0.5}{\input{tmp/bw_effort_v2.tex}}
	\begin{tablenotes}
		\scriptsize
		\item Note: farmer$\times$round panel; farmers report AEA visits and contacts. Columns (1)--(4):
		all 11 bi-weekly rounds (Apr--Sep 2025); columns (5)--(8): planting/establishment window only
		(rounds 05-1--06-2). Round and province FE; SE clustered at district. Control mean measured in
		the first round (all-rounds columns) or within the window. Wild-boot (Webb, 9{,}999) and RI
		(1{,}000) $p$-values reported. ***, **, * : 1, 5, 10\%.
	\end{tablenotes}
\end{frame}

\begin{frame}{Bi-weekly (farmer): Satisfaction to Extension Service}
	\centering
	\scalebox{0.46}{\input{tmp/bw_satisfy_v2.tex}}
	\begin{tablenotes}
		\scriptsize
		\item Note: farmer$\times$round panel. ``Satisfied'' = satisfied/very satisfied, defined for
		rounds with any visit; ``full'' recodes no-visit rounds to 0; ``add.\ ctrl'' adds visit and
		contact counts. Columns (1)--(4): all 11 rounds (Apr--Sep 2025); columns (5)--(8):
		planting/establishment window only (05-1--06-2), with the control mean measured within the
		window. Round and province FE; SE clustered at district. Wild-boot (Webb, 9{,}999) and RI
		(1{,}000) $p$-values. ***, **, * : 1, 5, 10\%.
	\end{tablenotes}
\end{frame}

\begin{frame}{Endline (AEA): Motivation and Locus of Control}
	\centering
	\scalebox{0.55}{\input{tmp/e_aea_attitude.tex}}
	\begin{tablenotes}
		\scriptsize
		\item Note: ANCOVA controlling for the baseline value; SE clustered at district ($N=43$ AEAs).
		Wild-boot = Webb wild cluster bootstrap (9{,}999); RI = randomization inference (1{,}000).
		***, **, * : 1, 5, 10\%.
	\end{tablenotes}
\end{frame}

\begin{frame}{Endline (AEA): Knowledge, Effort, Satisfaction}
	\centering
	\scalebox{0.58}{\input{tmp/e_aea_effort.tex}}
	\begin{tablenotes}
		\scriptsize
		\item Note: ANCOVA controlling for the baseline value; SE clustered at district ($N=43$).
		Self-reported visits winsorized at 1/99; three AEAs' reporting-period errors corrected in the
		revised data (Jul 2026), but the measure remains noisy --- the bi-weekly farmer reports are
		the preferred effort measure.
		***, **, * : 1, 5, 10\%.
	\end{tablenotes}
\end{frame}

\begin{frame}{Endline (Farmer): Production and Income}
	\centering
	\scalebox{0.46}{\input{tmp/e_farmer_prod.tex}}
	\begin{tablenotes}
		\scriptsize
		\item Note: ANCOVA with region FE and baseline-outcome control; SE clustered at district.
		Income = value of output at district-median prices $-$ cash costs (incl.\ hired labor);
		crop failure $=0$ output, so income can be negative. Log income uses the inverse
		hyperbolic sine. Wild-boot (Webb, 9{,}999) and RI (1{,}000) $p$-values reported. ***, **, * : 1, 5, 10\%.
	\end{tablenotes}
\end{frame}

\begin{frame}{Endline (Farmer): Sampling-Weighted Robustness}
	\centering
	\scalebox{0.46}{\input{tmp/e_farmer_weighted.tex}}
	\begin{tablenotes}
		\scriptsize
		\item Note: Population-representative estimates using inverse-probability sampling weights
		(rice-farmer population $/$ baseline sample size). Region FE and baseline-outcome control;
		SE clustered at district. ***, **, * : 1, 5, 10\%.
	\end{tablenotes}
\end{frame}

\begin{frame}{Endline (Farmer): Management-Practice Adoption}
	\centering
	\scalebox{0.5}{\input{tmp/e_farmer_practice.tex}}
	\begin{tablenotes}
		\scriptsize
		\item Note: ANCOVA with region FE and baseline-outcome control; SE clustered at district.
		***, **, * : 1, 5, 10\%.
	\end{tablenotes}
\end{frame}

\begin{frame}{Endline (Farmer): Sampling-Weighted Robustness (Practices)}
	\centering
	\scalebox{0.5}{\input{tmp/e_farmer_prac_w.tex}}
	\begin{tablenotes}
		\scriptsize
		\item Note: Population-representative estimates using inverse-probability sampling weights
		(rice-farmer population $/$ baseline sample size). Region FE and baseline-outcome control;
		SE clustered at district. ***, **, * : 1, 5, 10\%.
	\end{tablenotes}
\end{frame}

\begin{frame}{Endline: Key Takeaways}
\begin{itemize}\small
	\item \textbf{T1 (Feedback)} --- moved \emph{beliefs and service quality}:
	\begin{itemize}\small
		\item AEA: agricultural locus of control $+2.1$ (RI $p=0.001$); other attitudes unchanged
		\item Farmer: satisfaction $+9$--$18$pp (RI $p\approx0.02$); herbicide use and
		      crop failure fall; \emph{no} yield/income effect; \emph{no} effect on visit numbers
	\end{itemize}
	\item \textbf{T2 (Feedback + Training)} --- moved \emph{effort and knowledge}:
	\begin{itemize}\small
		\item AEA: general locus of control $+3.1$ (RI $p=0.01$); rice knowledge $+0.65$ (RI $p=0.05$)
		\item Farmer: visits $+0.15$/round ($\approx+31\%$ of control, RI $p=0.02$); satisfaction
		      $+6$--$11$pp (RI $p=0.02$); herbicide use falls; machine use up;
		      \emph{no} yield/income effect
	\end{itemize}
	\item \textbf{Stamp (GPS)} --- \emph{no systematic impact}: null on visits, satisfaction,
	      AEA attitudes/knowledge, and yield/income; no systematic practice effects
\end{itemize}
\end{frame}

% ======================= APPENDIX =======================
\appendix
\begin{frame}[plain]
	\vfill
	\centering
	{\Huge\bfseries\textcolor{myGreen}{Appendix}}
	\vfill
\end{frame}

% ---- A1. FB/TR reparameterization ----
\begin{frame}{Appendix: FB/TR --- Bi-weekly Effort}
	\centering
	\scalebox{0.5}{\input{tmp/bw_effort_fbtr.tex}}
	\begin{tablenotes}
		\scriptsize
		\item Note: same regression as the main table, reparameterized: FB $=1$\{T1 or T2\}
		(feedback effect $=$ T1 vs Control); TR $=1$\{T2\} (incremental effect of training on top
		of feedback $=$ T2 $-$ T1). Columns (1)--(4): all rounds; (5)--(8): planting window
		(05-1--06-2). Round and province FE; SE clustered at district; wild-boot (Webb, 9{,}999)
		and RI (1{,}000) $p$-values. ***, **, * : 1, 5, 10\%.
	\end{tablenotes}
\end{frame}

\begin{frame}{Appendix: FB/TR --- Bi-weekly Satisfaction}
	\centering
	\scalebox{0.46}{\input{tmp/bw_satisfy_fbtr.tex}}
	\begin{tablenotes}
		\scriptsize
		\item Note: FB $=1$\{T1 or T2\}; TR $=1$\{T2\} (T2 $-$ T1). Definitions as in the main
		satisfaction table. Columns (1)--(4): all rounds; (5)--(8): planting window (05-1--06-2).
		Round and province FE; SE clustered at district; wild-boot and RI $p$-values.
		***, **, * : 1, 5, 10\%.
	\end{tablenotes}
\end{frame}

\begin{frame}{Appendix: FB/TR --- AEA Attitudes (ANCOVA)}
	\centering
	\scalebox{0.62}{\input{tmp/e_aea_att_fbtr.tex}}
	\begin{tablenotes}
		\scriptsize
		\item Note: FB $=1$\{T1 or T2\}; TR $=1$\{T2\} (T2 $-$ T1). ANCOVA with baseline-outcome
		control; SE clustered at district ($N=43$); wild-boot and RI $p$-values. ***, **, * : 1, 5, 10\%.
	\end{tablenotes}
\end{frame}

\begin{frame}{Appendix: FB/TR --- AEA Knowledge, Effort, Satisfaction}
	\centering
	\scalebox{0.65}{\input{tmp/e_aea_eff_fbtr.tex}}
	\begin{tablenotes}
		\scriptsize
		\item Note: FB $=1$\{T1 or T2\}; TR $=1$\{T2\} (T2 $-$ T1). ANCOVA; SE clustered at
		district; wild-boot and RI $p$-values. ***, **, * : 1, 5, 10\%.
	\end{tablenotes}
\end{frame}

\begin{frame}{Appendix: FB/TR --- Farmer Production and Income}
	\centering
	\scalebox{0.46}{\input{tmp/e_farmer_prod_fbtr.tex}}
	\begin{tablenotes}
		\scriptsize
		\item Note: FB $=1$\{T1 or T2\}; TR $=1$\{T2\} (T2 $-$ T1). ANCOVA with region FE and
		baseline-outcome control; SE clustered at district; wild-boot and RI $p$-values.
		***, **, * : 1, 5, 10\%.
	\end{tablenotes}
\end{frame}

\begin{frame}{Appendix: FB/TR --- Farmer Management Practices}
	\centering
	\scalebox{0.5}{\input{tmp/e_farmer_prac_fbtr.tex}}
	\begin{tablenotes}
		\scriptsize
		\item Note: FB $=1$\{T1 or T2\}; TR $=1$\{T2\} (T2 $-$ T1). ANCOVA with region FE and
		baseline-outcome control; SE clustered at district; wild-boot and RI $p$-values.
		***, **, * : 1, 5, 10\%.
	\end{tablenotes}
\end{frame}

% ---- A2. Irrigation robustness (plot-level dummy) ----
\begin{frame}{Appendix: Plot-Level Irrigation Control (Farmer)}
	\centering
	\scalebox{0.46}{\input{tmp/e_farmer_prod_irrg.tex}}
	\begin{tablenotes}
		\scriptsize
		\item Note: replaces the scheme-level design dummy with the \emph{plot-level} irrigation
		indicator (QY1.1\_9). The plot-level indicator is a farmer choice and potentially
		post-treatment (``bad control''): read descriptively. ANCOVA with region FE and
		baseline-outcome control; SE clustered at district; wild-boot and RI $p$-values.
		***, **, * : 1, 5, 10\%.
	\end{tablenotes}
\end{frame}

% ---- A3. 2SLS (4 endogenous x production/adoption, two panels per slide) ----
\begin{frame}{Appendix: 2SLS --- Extension Intensity and Production}
	\centering
	{\footnotesize\textbf{Endogenous: cumulative visits (Apr--Sep)}}\\[1pt]
	\scalebox{0.46}{\input{tmp/e_2sls_o1_cumQ1.tex}}\\[4pt]
	{\footnotesize\textbf{Endogenous: cumulative contacts (Apr--Sep)}}\\[1pt]
	\scalebox{0.46}{\input{tmp/e_2sls_o1_cumQ2.tex}}
	\begin{tablenotes}
		\scriptsize
		\item Note: 2SLS; endogenous regressor $=$ cumulative farmer-reported visits (top) or
		contacts (bottom) over the 11 bi-weekly rounds, instrumented by T1/T2 assignment.
		ANCOVA controls as in the main tables; SE clustered at district. KP $F$ = Kleibergen--Paap
		weak-instrument $F$ ($\approx$10 clusters: weak-IV caution); Hansen $J$ from the
		het.-robust VCE. ***, **, * : 1, 5, 10\%.
	\end{tablenotes}
\end{frame}

\begin{frame}{Appendix: 2SLS --- Satisfaction and Production}
	\centering
	{\footnotesize\textbf{Endogenous: mean satisfaction level (Apr--Sep), visited rounds}}\\[1pt]
	\scalebox{0.46}{\input{tmp/e_2sls_o1_msatall.tex}}\\[4pt]
	{\footnotesize\textbf{Endogenous: mean satisfaction level (05-1--06-2), visited rounds}}\\[1pt]
	\scalebox{0.46}{\input{tmp/e_2sls_o1_msat.tex}}
	\begin{tablenotes}
		\scriptsize
		\item Note: endogenous regressor $=$ mean 1--5 satisfaction level over visited rounds,
		full period (top) or May--June only (bottom); the May--June measure \emph{predates
		harvest}, limiting reverse causality from realized outcomes. Farmers with $\geq$1 contact
		in the respective period. Instruments, controls, and statistics as on the previous slide.
		***, **, * : 1, 5, 10\%.
	\end{tablenotes}
\end{frame}

\begin{frame}{Appendix: 2SLS --- Extension Intensity and Technology Adoption}
	\centering
	{\footnotesize\textbf{Endogenous: cumulative visits (Apr--Sep)}}\\[1pt]
	\scalebox{0.48}{\input{tmp/e_2sls_o2_cumQ1.tex}}\\[4pt]
	{\footnotesize\textbf{Endogenous: cumulative contacts (Apr--Sep)}}\\[1pt]
	\scalebox{0.48}{\input{tmp/e_2sls_o2_cumQ2.tex}}
	\begin{tablenotes}
		\scriptsize
		\item Note: as the first 2SLS slide, for the technology-adoption outcomes.
		***, **, * : 1, 5, 10\%.
	\end{tablenotes}
\end{frame}

\begin{frame}{Appendix: 2SLS --- Satisfaction and Technology Adoption}
	\centering
	{\footnotesize\textbf{Endogenous: mean satisfaction level (Apr--Sep), visited rounds}}\\[1pt]
	\scalebox{0.48}{\input{tmp/e_2sls_o2_msatall.tex}}\\[4pt]
	{\footnotesize\textbf{Endogenous: mean satisfaction level (05-1--06-2), visited rounds}}\\[1pt]
	\scalebox{0.48}{\input{tmp/e_2sls_o2_msat.tex}}
	\begin{tablenotes}
		\scriptsize
		\item Note: as the second 2SLS slide, for the technology-adoption outcomes.
		***, **, * : 1, 5, 10\%.
	\end{tablenotes}
\end{frame}

\end{document}
